\chapter{Linear Algebra}

% \red{The following is from Math-tools, a M.Sc class}
\subsection{recap}
\begin{definition}
A vector space over field $\mathbb{F}$ is a set $V$ with two operations:
Vector addition and scalar multiplication.
\end{definition}

\begin{definition}
$T:V\to W$ ($W,V$ are linear spaces) is a linear map ``if it respects''
the space's linear structure. That is $T\left(u+v\right)=T\left(u\right)+t\left(v\right)$
and $\forall\alpha$, $T\left(\alpha u\right)=\alpha T\left(u\right)$.
\end{definition}

\begin{definition}
Let $V$ be a f.d (finite dimension) vector space $\Rightarrow$ there
are $v_{1}\dots v_{d}\in V$ s.t. $\forall v\in V\exists\alpha_{1}\dots\alpha_{d}$
s.t. $v=\sum\alpha_{i}v_{i}$. The smallest such $d$ is called $dim\left(V\right)$,
and a set $v_{1}\dots v_{d}$ -- a basis.
\end{definition}

\begin{claim}
Let us fix a basis. Then every $v\in V$ is uniquely determined by
the scalars $\alpha_{1}\dots\alpha_{d}$ 
\begin{align*}
v & \leftrightarrow\left(\alpha_{1},\dots\alpha_{n}\right)\\
v_{1} & \leftrightarrow\left(1,0,\dots0\right)
\end{align*}
\end{claim}

\begin{claim}
Every linear mapping $T:V\to W$ ($W$ has a standard basis) is represented
by a matrix $A_{k\times d}$ ($dim\left(W\right)=k$.
\end{claim}

\[
T\left(v\right)=T\left(\alpha_{1},\dots,\alpha_{n}\right)=A\alpha
\]
 %
\begin{insight}
what? $A\alpha$?

What are infinite dim vector spaces? Consider the space of all continuous
functions $\left[a,b\right]\to\R$.%
\end{insight}

\begin{definition}
A vector $v$ is an \textbf{eigenvector} of linear map $T$ with \textbf{eigenvalue}
$\lambda$ of $T$ if $T\left(v\right)=\lambda\cdot v$.
\end{definition}


\subsubsection{\textquotedblleft Rough statements\textquotedblright{} }
\begin{itemize}
\item Linear maps are comprised of rotations and stretching.
\item ``Size of a vector'': The most important example: $L_{2}$ norm.
The size of the vector $\left|\left|\left(x_{1},\dots x_{n}\right)\right|\right|_{2}=\sqrt{\sum x_{i}^{2}}$
.
\item Rotation: a linear map that preserves the $L_{2}$ norm, $\forall v,\left|\left|T\left(v\right)\right|\right|_{2}=\left|\left|v\right|\right|_{2}$.
\item When do we say that $T$ stretches the vector $v$ by $\lambda$?
$T\left(v\right)=\lambda v$. The more standard term is: $v$ is an
eigenvector of $T$ with eigenvalue of $\lambda$.
\end{itemize}
Let's make things more concrete.

\section{Metric spaces}
\begin{definition}
A metric space is a pair $\left(X,d\right)$ where $X$ is a set and
$d$ is a distance function $d:X\times X\to\R$ such that:
\begin{enumerate}
\item $d\left(x,y\right)=d\left(y,x\right)$
\item $\forall x,y\in X$ $d\left(x,y\right)\ge0$ and equality to zero
iff $x=y$. 
\item triangle inequality -- $\forall x,y,z\in X$ , $d\left(x,y\right)+d\left(y,z\right)\ge d\left(x,z\right)$
\end{enumerate}
\end{definition}


\section{normed spaces}

``From size to distance'' (in linear spaces) $\left|\left|u\right|\right|$
``the size of the vector'' we can define as $d\left(u,v\right):=\left|\left|u-v\right|\right|$%
\begin{insight}
like in real numbers, the size is it's abs value, and the distance
is the subtract of two numberss inside abs value%
\end{insight}
.
\begin{definition}
Let $V$ be a real linear space, $\left|\left|\cdot\right|\right|:V\to R$
is a norm if for every $u,v\in V$ and $\alpha\in\R$ it has the following
properties:
\begin{enumerate}
\item $\forall v\in V,\left|\left|v\right|\right|\ge0,\left|\left|v\right|\right|=0\iff v=0$
\item $\forall\alpha\in\R,v\in V,$ $\left|\left|\alpha v\right|\right|=\left|\alpha\right|\cdot\left|\left|v\right|\right|$
\item triangle inequality -- $\forall u,v\in V,\left|\left|u+v\right|\right|\le\left|\left|u\right|\right|+\left|\left|v\right|\right|$
\end{enumerate}
\end{definition}


\subsection{Important examples of norms:}

$L_{p}$ norms, where $1\le p\le\infty$ are defined as follows: $\left|\left|\left(x_{1},\dots,x_{n}\right)\right|\right|_{p}:=\left(\sum\left|x_{i}\right|^{p}\right)^{\frac{1}{p}}$

The most important norms are for $p=1,2,\infty$
\begin{enumerate}
\item $\left|\left|\left(x_{1},\dots,x_{n}\right)\right|\right|_{\infty}:=\max\left\{ \left|x_{i}\right|\right\} $
\item $\left|\left|\left(x_{1},\dots,x_{n}\right)\right|\right|_{1}:=\sum\left|x_{i}\right|$
(Manhattan distance)
\item $\left|\left|\left(x_{1},\dots,x_{n}\right)\right|\right|_{2}:=\left(\sum\left|x_{2}\right|^{2}\right)^{\frac{1}{2}}$
\end{enumerate}
The fact that $L_{p}$ is a norm is not trivial. (The triangle inequality
requires proof)

\subsection{Balls}

Associated with every norm $\left|\left|\cdot\right|\right|$ is its
\uline{unit ball} on a space $V$: $B_{\left|\left|\cdot\right|\right|}:=\left\{ x\in V:\left|\left|x\right|\right|\le1\right\} $.
\begin{itemize}
\item The unit ball of $L_{2}$ is a round ball. 
\item The unit ball for $L_{\infty}$ is a cube: $\left|\left|\left(x_{1},\dots x_{n}\right)\right|\right|_{\infty}\le1\iff\forall i,-1\le x_{i}\le1$.
\item The unit ball of $L_{1}$ in 3 dims is an octahedron ($\left|x\right|+\left|y\right|+\left|z\right|\le1$
in $\R^{3}$ ).
\end{itemize}

\subsubsection{inequalities}

how many vertices, edges and inequalities do we have?
\begin{itemize}
\item cube: 
\begin{align*}
-1 & \le x_{1}\le1\\
-1 & \le x_{2}\le1\\
 & \vdots\\
-1 & \le x_{n}\le1
\end{align*}
 So we have $2^{n}$ vertices.
\item Cross polytope: (the $L_{1}$): $\pm e_{1},\pm e_{2},\dots$ In $\R^{3}$:
\begin{align*}
x+y+z & \le1\\
-x+y+z & \le1\\
x-y+z & \le1\\
 & \vdots
\end{align*}
 we have 8 inequalities in $\R^{3}$ and thus $8$ facets.
\end{itemize}
\begin{fact}
$\forall$ norm $\left|\left|\cdot\right|\right|$, its unit ball
$\ball{}$ is :
\begin{enumerate}
\item centrally symmetric. ($x\in B\Rightarrow-x\in B$ $\left|\left|-x\right|\right|=\left|-1\right|\cdot\left|\left|x\right|\right|\le1$
)
\item convex %
\begin{insight}
if $x_{1}$ and $x_{2}$ are in $B$ then the interval defined by
$\left\{ \lambda x_{1}+\left(1-\lambda\right)x_{2}:\lambda\in\left[0,1\right]\right\} \subseteq B$.%
\end{insight}
\item full dimensional set in $\R^{n}$ %
\begin{insight}
That is if $x\in B$ then $-x\in B$.%
\end{insight}
\end{enumerate}
\end{fact}

\begin{proposition}
if $B\subseteq\R^{n}$ and $B$ has the properties: 
\begin{enumerate}
\item centrally symmetric. ($x\in B\Rightarrow-x\in B$ $\left|\left|-x\right|\right|=\left|-1\right|\cdot\left|\left|x\right|\right|\le1$
)
\item convex %
\begin{insight}
if $x_{1}$ and $x_{2}$ are in $B$ then the interval defined by
$\left\{ \lambda x_{1}+\left(1-\lambda\right)x_{2}:\lambda\in\left[0,1\right]\right\} \subseteq B$.%
\end{insight}
\item full dimensional set in $\R^{n}$ %
\begin{insight}
That is if $x\in B$ then $-x\in B$.%
\end{insight}
{} 
\end{enumerate}
then $\exists\left|\left|\cdot\right|\right|$ s.t $B=B_{\left|\left|\cdot\right|\right|}$.

\end{proposition}

\begin{definition}[isometry]
 Let $\left(X,\left|\left|\cdot\right|\right|\right)$ be a normed
space, a map $T:X\to X$ is an isometry if $\forall x\left|\left|Tx\right|\right|=\left|\left|x\right|\right|$
.
\end{definition}

Another key property of $L_{2}$ is that this norm is derived from
an \uline{inner product}. $\left\langle x,y\right\rangle =\sum x_{i}y_{i}$.
and $\left|\left|x\right|\right|=$$\left\langle x,x\right\rangle ^{\frac{1}{2}}$.

We associate with every $x\in\R^{n}$ the map $y\to\left\langle x,y\right\rangle $
. the map: $\left\langle x,\cdot\right\rangle \left(y\right)$
\begin{definition}[dual norm]
 Let $\left|\left|\cdot\right|\right|$ be a norm on $\R^{n}$, it's
dual norm is $\left|\left|x\right|\right|^{*}=\sup\limits _{y\in\R^{n}}\frac{\left|\left\langle x,y\right\rangle \right|}{\left|\left|y\right|\right|}=\sup\limits _{\left|\left|y\right|\right|=1}\left|\left\langle x,y\right\rangle \right|$%
\begin{insight}
We can write $\max$ instead of $\sup$%
\end{insight}
{} %
\begin{insight}
it's equal because $y\cdot\alpha$ in$\frac{\left|\left\langle x,\alpha y\right\rangle \right|}{\left|\left|\alpha y\right|\right|}=\frac{\left|\left\langle x,y\right\rangle \right|}{\left|\left|y\right|\right|}$
so the expression doesn\'{t} change for any $\alpha$ so we can multiplie
any $y$ by some number to get $\left|\left|y\right|\right|=1$ and
the norm stays the same%
\end{insight}
.
\end{definition}

\begin{xca}
What is the dual norm of $L_{\infty}$?
\begin{sol}
$x\in\R^{n}$ ,
\[
\left|\left|x\right|\right|_{\infty}^{*}=\max_{\left|\left|y\right|\right|_{\infty}=1}\left|\left\langle x,y\right\rangle \right|=\left|\left|x\right|\right|_{1}
\]

\begin{insight}
to start it up lets use some real vector to get a hunch : $x=\left(7,-3,5,9\right)$
so we want $y$ with these signs to the coordinates. $y=\left(+,-,+,+\right)$
so the best $y$ is such that $\left|\left|y\right|\right|_{\infty}=1$
is $y=\left(1,-1,1,1\right)$. %
\begin{insight}
$L_{1}$ and $L_{\infty}$ are the dual norms of each other.%
\end{insight}
\end{insight}

\end{sol}

\end{xca}

\begin{example}
What is the dual norm of $L_{2}$? $x\in\R^{n}$. what is $\left|\left|x\right|\right|_{2}^{*}$?
\[
\left|\left|x\right|\right|_{2}^{*}=\max_{\left|\left|y\right|\right|_{2}=1}\left|\left\langle x,y\right\rangle \right|
\]
 By Cauchy Schwartz the best choice is $y=\frac{x}{\left|\left|x\right|\right|_{2}}$
\begin{insight}
$\left\langle x,y\right\rangle =\left|\left|x\right|\right|\cdot\left|\left|y\right|\right|\cdot cos\left(\alpha\right)$
where $\alpha$ is the angle between the two, thus we want $\cos\alpha=1$
when $x=y$ else this is smaller then it.%
\end{insight}
\[
\left|\left|x\right|\right|_{2}^{*}=\max_{\left|\left|y\right|\right|_{2}=1}\left|\left\langle x,y\right\rangle \right|=\left\langle x,\frac{x}{\left|\left|x\right|\right|_{2}}\right\rangle =\frac{\left(\left|\left|x\right|\right|_{2}\right)^{2}}{\left|\left|x\right|\right|_{2}}=\left|\left|x\right|\right|_{2}
\]
\end{example}

\begin{claim}
If $1\le p,q\le\infty$, $\frac{1}{q}+\frac{1}{p}=1\Rightarrow\left|\left|x\right|\right|_{p}^{*}=\left|\left|x\right|\right|_{q}$.
\end{claim}

\begin{insight}
Find $x$ such that $Ax\thickapprox b$. no such $x$ (too many constraints)
Find $x$ such that $\left|\left|Ax-b\right|\right|$ is minimized%
\end{insight}

\begin{insight}

\subsection{some reminders from previous class(?)}

Norm on real vector space $V$ is a map $\left|\left|\cdot\right|\right|:V\to\R$
if
\begin{enumerate}
\item $\forall x\in V$ $\left|\left|x\right|\right|\ge0$ and there\textquoteright s
equality if and only if $x=0$
\item Linearity: $\forall v\in V$$\lambda\in\R,\left|\left|\lambda v\right|\right|=\left|\lambda\right|\cdot\left|\left|v\right|\right|$
\item triangle inequality: $\left|\left|x+y\right|\right|\le\left|\left|x\right|\right|+\left|\left|y\right|\right|$
\end{enumerate}
%

\subsubsection{The geometric perspective:}

corresponding to $\left|\left|\cdot\right|\right|$ is it's unit ball
$B_{\left|\left|\cdot\right|\right|}=\left\{ x\in V|\left|\left|x\right|\right|\le1\right\} $
\begin{fact}
B is
\begin{enumerate}
\item convex %
\begin{insight}
(coming from traingle inequallity in norms)%
\end{insight}
\item centrally symmetric
\item full dimensional subset of $V$. Means ``it contains a Euclidean''
\begin{insight}
any vector who is not zero has positive norm%
\end{insight}
\end{enumerate}
\end{fact}

on the other hand: every set $B\subseteq V$ with propeties $\left(1\right),\left(2\right),\left(3\right)$
is the unit ball of some norm.
\begin{example}
Very important examples for norms are: $l_{p}$%
\begin{insight}
\"{s}mall $l_{p}$ is when we work with vectors (for sequnce spaces)
and $L_{p}$ is for function spaces%
\end{insight}
{} norm on $\R^{n}$ $\left|\left|x\right|\right|_{p}:=\left(\sum\left|x_{i}\right|^{p}\right)^{\frac{1}{p}}$
in perticular $p=1,2,3,\dots\infty$
\end{example}

\end{insight}

\begin{claim}
For $\infty\ge p\ge1$, $l_{p}$ is a norm.

We prove only the triangle inequality because the other values are
obvious. We need to show: $x,y\in\R^{n}$ $\left|\left|x+y\right|\right|_{p}\le\left|\left|x\right|\right|_{p}+\left|\left|y\right|\right|_{p}$.
We start with Young's inequality
\begin{proposition}
Let $p,q>0,\frac{1}{p}+\frac{1}{q}=1$ , $a,b>0$ then $a\cdot b\le\frac{1}{p}a^{p}+\frac{1}{q}b^{q}$.
\begin{proof}
recall that ``$log$'' is a concave function. %
\begin{insight}
take any two points $x,y$ on the curve of the function, they are
below the curve of the log ($\forall x,y$ $1\ge\lambda\ge0$, $\lambda\log\left(x\right)+\left(1-\lambda\right)\log y\le\log\left(\lambda x+\left(1-\lambda\right)y\right)$)%
\end{insight}
{} then 
\begin{align*}
\log ab & =\log\left(a\right)+\log b=\frac{1}{p}\cdot p\log\left(a\right)+\frac{1}{q}\cdot q\log b\\
 & =\frac{1}{p}\log\left(a^{p}\right)+\frac{1}{q}\log\left(b^{q}\right)\le\text{\ensuremath{\log\left(\frac{1}{p}a^{p}+\frac{1}{q}b^{q}\right)}}
\end{align*}
\end{proof}
\begin{proposition}[Holder's inequality]
$\forall x,y\in\R^{n}$ $\left|\left\langle x,y\right\rangle \right|\le\left|\left|x\right|\right|_{p}\left|\left|y\right|\right|_{q}$
if $\frac{1}{p}+\frac{1}{q}=1$ $p,q>0$.
\begin{proof}
WLOG (without loss of generality) $x,y\ge0$. $\sum x_{j}y_{j}\overset{?}{\le}\left(\sum\left|x_{i}\right|^{p}\right)^{\frac{1}{p}}\cdot\left(\sum\left|x_{i}\right|^{q}\right)^{\frac{1}{q}}$.
\begin{align*}
\frac{\left\langle x,y\right\rangle }{\left|\left|x\right|\right|_{p}\left|\left|y\right|\right|_{q}} & =\sum\limits _{j}\frac{x_{j}}{\left|\left|x\right|\right|_{p}}\frac{y_{j}}{\left|\left|x\right|\right|_{q}}\le\sum\left[\frac{1}{p}\left(\frac{x_{j}}{\left|\left|x\right|\right|_{p}}\right)^{p}+\frac{1}{q}\left(\frac{y_{j}}{\left|\left|y\right|\right|_{y}}\right)^{q}\right]=\\
 & =\frac{1}{p}\frac{\sum x_{j}^{p}}{\left|\left|x\right|\right|_{p}^{p}}+\frac{1}{q}\frac{\sum y_{j}^{q}}{\left|\left|y\right|\right|_{q}^{q}}=\frac{1}{p}+\frac{1}{q}=1
\end{align*}
 we can multiply the denominator and get what we wanted.
\end{proof}
\end{proposition}

\end{proposition}

\begin{proof}
What we want to show $\forall x,y\in\R^{n}$ $\left(\sum\left|x_{j}+y_{j}\right|^{p}\right)^{\frac{1}{p}}\le\left(\sum\left|x_{j}\right|^{p}\right)^{\frac{1}{p}}+\left(\sum\left|y_{j}\right|^{p}\right)^{\frac{1}{p}}$.
$l_{p}$ is a norm: again WLOG $x,y\ge0$ then 
\begin{align*}
\left|\left|x+y\right|\right|_{p}^{p} & =\sum_{j}\left(x_{j}+y_{j}\right)^{p}=\sum\left(x_{j}+y_{j}\right)\left(x_{j}+y_{j}\right)^{p-1}=\\
 & =\sum x_{j}\left(x_{j}+y_{j}\right)^{p-1}+y_{j}\left(x_{j}+y_{j}\right)^{p-1}=\\
 & =\left[\sum x_{j}\underbrace{\left(x_{j}+y_{j}\right)^{p-1}}_{v_{j}}\right]+\left[\sum y_{j}\underbrace{\left(x_{j}+y_{j}\right)^{p-1}}_{vj}\right]
\end{align*}
 (we will talk about $v_{j}$ in a moment) so %
\begin{insight}
we use holder inequallity with $\left|\left\langle u,v\right\rangle \right|\le\left|\left|u\right|\right|_{p}\left|\left|v\right|\right|_{q}$%
\end{insight}
{} with holders: 
\begin{align*}
 & \le\left|\left|x\right|\right|_{p}\left|\left|x+y\right|\right|_{p}^{p-1}+\left|\left|y\right|\right|_{p}\cdot\left|\left|x+y\right|\right|_{p}^{p-1}=\\
 & \left|\left|x+y\right|\right|_{p}^{p-1}\left(\left|\left|x\right|\right|_{p}+\left|\left|y\right|\right|_{p}\right)
\end{align*}
divide by $\left|\left|x+y\right|\right|_{p}^{p-1}$ and we get what
we wanted.\\
Lets get back to the missing part: If $v_{j}=\left(x_{j}+y_{j}\right)^{p-1}\Rightarrow\left|\left|v\right|\right|_{q}=\left|\left|x+y\right|\right|_{p}^{p-1}$
explanation: $\frac{1}{q}+\frac{1}{p}=1\iff q=\frac{p}{p-1}$ 
\begin{align*}
\left|\left|v\right|\right|_{q} & =\left(\sum v_{j}^{q}\right)^{\frac{1}{q}}=\left(\sum\left(x_{j}+y_{j}\right)^{q\left(p-1\right)}\right)^{\frac{1}{q}}=\left(\sum\left(x_{j}+y_{j}\right)^{p}\right)^{\frac{1}{q}}=\\
 & =\left(\sum\left(x_{j}+y_{j}\right)^{p}\right)^{\frac{p-1}{p}}=\left(\left(\sum\left(x_{j}+y_{j}\right)^{p}\right)^{\frac{1}{p}}\right)^{p-1}=\left|\left|x+y\right|\right|_{p}^{p-1}
\end{align*}
 Now using this $v_{j}$ we can simply use holder's inequality to
get $\le\left|\left|x\right|\right|_{p}\left|\left|x+y\right|\right|_{p}^{p-1}+\left|\left|y\right|\right|_{p}\cdot\left|\left|x+y\right|\right|_{p}^{p-1}=$
and finish our proof.
\end{proof}
\end{claim}

\begin{definition}[operator norm]
 $\left(U,\left|\left|\cdot\right|\right|_{U}\right)\overset{\psi}{\to}\left(V,\left|\left|\cdot\right|\right|_{v}\right)$
. 
\[
\left|\left|\psi\right|\right|:=\sup_{u\in U}\frac{\left|\left|\psi u\right|\right|_{V}}{\left|\left|u\right|\right|_{U}}
\]
 If $\psi$ is linear then 
\[
\left|\left|\psi\right|\right|:=\sup_{u\in U}\frac{\left|\left|\psi u\right|\right|_{V}}{\left|\left|u\right|\right|_{U}}=\sup_{\begin{array}{c}
u\in U\\
\left|\left|u\right|\right|_{U}=1
\end{array}}\left|\left|\psi u\right|\right|_{V}
\]
 
\end{definition}

an important special case: $V=\R,\left|\left|\cdot\right|\right|_{V}=\left|\cdot\right|$.
to every $x\in\R^{n}$ corresponds a linear map $\R^{n}\to\R,y\in\R^{n}\to\left\langle x,y\right\rangle $,
``We think of x as a linear functional''. So the norm of the operator
$\left\langle x,\cdot\right\rangle $: $\left|\left|x\right|\right|^{*}:=\max\limits _{\left|\left|y\right|\right|=1}\left|\left\langle x,y\right\rangle \right|$.%
\begin{insight}
Therefore if $\left|\left|\cdot\right|\right|$ is a norm on $\R^{n}$
this induces a new norm called the dual of $\left|\left|\cdot\right|\right|$
denoted $\left|\left|x\right|\right|^{*}:=\max\limits _{\left|\left|y\right|\right|=1}\left|\left\langle x,y\right\rangle \right|$.%
\end{insight}
\begin{insight}
to understand the notations: $\left|\left|x\right|\right|^{*}:\left|\left|\left\langle x,\cdot\right\rangle \right|\right|\max\limits _{\left|\left|y\right|\right|=1}\left|\left\langle x,y\right\rangle \right|$%
\end{insight}
{} 
\begin{definition}
The unit ball $B^{*}$ of the dual norm is called the polar body of
$B$ and is defined as follows:

\[
B^{*}=\left\{ y\in\R^{n}|\forall x\in B\left\langle x,y\right\rangle \le1\right\} 
\]
\end{definition}

\begin{example}
Question: Let $B$ be the unit ball of some norm $\left|\left|\cdot\right|\right|$.
What can we say about the unit ball of the dual norm?
\begin{itemize}
\item $l_{\infty}$ Ball is $\left[-1,1\right]^{n}$ $2^{n}$ vertices $2n$
facets.
\item $l_{1}$ Ball is the cross polytope $\sum\left|x_{i}\right|\le1$.
Has $2n$ vertices (the standard basis $\pm e_{i}$), and $2^{n}$
facets $\epsilon\in\left\{ -1,1\right\} ^{n}$ $\sum\epsilon_{i}\cdot x_{i}\le1$.
\end{itemize}
What is the polar $B^{*}$ of the unit cube $B=\left[-1,1\right]^{n}$?
\[
B^{*}=\left\{ y\in\R^{n}|\begin{array}{c}
-1\le x_{1}\le1,-1\le x_{2}\le1\dots\\
\sum x_{i}y_{i}\le1
\end{array}\right\} 
\]
 So given $y$ the ``most difficult'' $x$ to satisfy %
\begin{insight}
$x$ that maximizes this expression%
\end{insight}
. $x_{i}=\text{sign}\left(y_{i}\right)$ in which case $\sum x_{i}y_{i}=\sum\left|y_{i}\right|$.
\end{example}


\section{\textquotedblleft What do linear maps do\textquotedblright ?}

$A$ is a matrix

\subsection{Stretch}

If $Av=\lambda v$ holds we say that $v$ is an eigenvector of $A$
with eigenvalue $\lambda$. \\
$\left(A-\lambda I\right)v=0$ i.e. $v$ is in the kernel of $\left(A-\lambda I\right)$.
In particular $\left(A-\lambda I\right)$ is \textbf{singular}, hence
$\det\left(A-\lambda I\right)=0$, and because $\det\left(A-\lambda I\right)$
is polynomial in $\lambda$. ``The characteristic polynomial $p_{A}$
of $A$'' $\Rightarrow$ Every eigenvalue is a root of $p_{A}$.
\begin{proposition}
If A is a real $n\times n$ matrix and if $u_{1}\dots u_{n}$ are
eigenvectors with \uline{distinct} eigenvalues then $u_{1}\dots u_{n}$
are linearly independent.

\begin{insight}
This is useful in case $u_{i}$ forms a basis. $x=\sum\alpha_{i}u_{i}$
then $Ax=\sum\alpha_{i}\lambda_{i}u_{i}$ then we are in a good situation.
``We have good understanding of the linear operator''%
\end{insight}
\begin{proof}
Assume to the contrary that they are linearly dependent and let $k$
be the smallest size of a linearly dependent subset. 
\[
\sum\limits _{i=1}^{k}b_{i}u_{i}=0\Rightarrow A\left(\sum_{i=1}^{k}b_{i}u_{i}\right)=\left(\sum_{i=1}^{k}b_{i}\lambda_{i}u_{i}\right)=0
\]
furthermore we know that $\sum_{i=1}^{k}\left(b_{i}u_{i}\right)=0$
(from our assumption) so
\[
0=\sum_{i=1}^{k}\left(b_{i}u_{i}\right)=\lambda_{k}\sum_{i=1}^{k}\left(b_{i}u_{i}\right)=\sum_{i=1}^{k}\left(b_{i}\lambda_{k}u_{i}\right)
\]
 Let us subtract the first from the second:
\begin{align*}
0 & =\sum_{i=1}^{k}b_{i}\left(\lambda_{i}-\lambda_{k}\right)u_{i}=\\
 & =b_{1}\left(\lambda_{1}-\lambda_{k}\right)u_{1}+\dots+b_{k-1}\left(\lambda_{k-1}-\lambda_{k}\right)u_{k-1}+b_{k}\left(\lambda_{k}-\lambda_{k}\right)u_{k}=\\
 & =b_{1}\left(\lambda_{1}-\lambda_{k}\right)u_{1}+\dots+b_{k-1}\left(\lambda_{k-1}-\lambda_{k}\right)u_{k-1}+0=\\
 & =\sum_{i=1}^{k-1}b_{i}\left(\lambda_{i}-\lambda_{k}\right)u_{i}
\end{align*}
 We found a shorter non trivial linear dependency, which is a contradiction. 
\end{proof}
How should we deal with general linear transformation? we should find
eigenvectors and if they are a basis we can use it to understand the
matrix. but what do we do when this isn't the case? What can go wrong? 
\end{proposition}

\begin{example}
example of a matrix where the eigenvectors do not form a basis:
\end{example}

\[
A=\left(\begin{array}{cccc}
0 & 1 & \dots & 0\\
0 & \ddots & \ddots & \vdots\\
\vdots & \ddots & \ddots & 1\\
0 & \dots & 0 & 0
\end{array}\right)
\]
 
\[
\det\left(\lambda I-A\right)=\left|\begin{array}{cccc}
\lambda & -1 & \dots & 0\\
0 & \ddots & \ddots & \vdots\\
\vdots & \ddots & \ddots & -1\\
0 & \dots & 0 & \lambda
\end{array}\right|=\lambda^{n}=0
\]
$\lambda=0$ is the only eigenvalue.
\[
A\left(\begin{array}{c}
x_{1}\\
\vdots\\
x_{n-1}\\
x_{n}
\end{array}\right)=\left(\begin{array}{c}
x_{2}\\
\vdots\\
x_{n}\\
0
\end{array}\right)=0\cdot x
\]
 $x_{2}=x_{3}=\dots=x_{n}=0$ the set of eigenvectors corresponding
to $\lambda=0$ is $\left\langle e_{1}\right\rangle $. $A^{n}=0$
(When we reapply $A$ it maps to zero after enough times) How to deal
with it, at least in practical situations?
\begin{definition}
Say that a matrix $B$ $n\times n$ is ``problem free'' if it has
distinct eigenvalues.
\end{definition}

\begin{fact}
``Most'' matrices are problem free. %
\begin{insight}
thats why we don\'{t} need to deal with the ones that don't%
\end{insight}
\end{fact}


\subsection{Rotations}

Linear maps which are isometries. It suffices that $l_{2}$ norms
are preserved $\left|\left|x-y\right|\right|_{2}^{2}=\left|\left|x\right|\right|_{2}^{2}+\left|\left|y\right|\right|_{2}^{2}-2\left\langle x,y\right\rangle $
\begin{insight}
the only norm for which such things exist is $l_{2}$.%
\end{insight}

\begin{definition}
A real matrix $A$ is called orthogonal if $A^{T}A=I$.

\begin{insight}
(This means that they are inverses too)%
\end{insight}

\end{definition}

\begin{definition}
A complex matrix $U$ is called unitary if $UU^{*}=I$ %
\begin{insight}
where $U^{*}=\overline{U^{T}}$?%
\end{insight}
\end{definition}

\begin{claim}
$\left\langle x,Ay\right\rangle =\left\langle A^{T}x,y\right\rangle $
\begin{proof}
in general $\left\langle u,v\right\rangle =\sum u_{i}v_{i}$, so
\begin{align*}
\left\langle x,Ay\right\rangle  & =\sum x_{i}\left(Ay\right)_{i}=\sum_{i}x_{i}\sum_{j}a_{ij}y_{j}=\sum_{i,j}a_{ij}x_{i}y_{j}=\\
\end{align*}
 The $a_{ij}$ of $A^{T}$.
\[
=\sum_{i,j}a_{ji}^{T}x_{i}y_{j}=\sum_{j}y_{j}\sum_{i}a_{ji}^{T}x_{i}=\left\langle A^{T}x,y\right\rangle 
\]
\begin{insight}
$a_{ji}^{T}$ is like take the original matrix\c{ }transpose it and
then look at $a_{ji}$%
\end{insight}
\end{proof}
\end{claim}

\begin{thm}
The linear map $\R^{n}\to\R^{n}$ corresponding to a matrix $A$ is
an isometry if and only if $A$ is orthogonal.
\begin{proof}
Suppose that $A$ is orthogonal. We want to show that $\forall x$
$\left|\left|Ax\right|\right|_{2}=\left|\left|x\right|\right|_{2}$.
need to show :
\[
\left\langle Ax,Ax\right\rangle \overset{?}{=}\left\langle x,x\right\rangle 
\]
we know that $\left\langle Ax,By\right\rangle =\left\langle B^{\top}Ax,y\right\rangle $
so

\[
\left\langle Ax,Ax\right\rangle =\left\langle A^{\top}Ax,x\right\rangle =\left\langle Ix,x\right\rangle =\left\langle x,x\right\rangle 
\]
 as wanted.
\end{proof}
\end{thm}

\begin{claim}
An orthogonal matrix $A$ is one that satisfies $AA^{T}=I$($\iff A^{T}A=I$)
\begin{proof}
the proof is for $AA^{T}=I\iff A^{T}A=I$: $A$ has a right inverse,
(namely $A^{T}$) $\Rightarrow$ it has a left inverse, say $B$ $BA=I$
\[
A^{T}=IA^{T}=BAA^{T}=BI=B
\]
 and we have what we wanted.
\end{proof}
\end{claim}

.What does $\left(HH^{T}\right)_{i,j}$ equals to? $\left(HH^{T}\right)_{i,j}=$
the inner product of the $i$th row of $H$ and it's $j$th row. $H^{T}H=$
the inner product of the $i$th column of $H$ and the $j$th column.
\begin{proposition}
A real square matrix $A$ is orthogonal iff $x\to Ax$ is an isometry,
i.e. $\forall x\left|\left|Ax\right|\right|_{2}=\left|\left|x\right|\right|$.
\begin{proof}
If $A$ is orthogonal: $\left|\left|Ax\right|\right|^{2}=\left\langle Ax,Ax\right\rangle =\left\langle A^{T}Ax,x\right\rangle =\left|\left|x\right|\right|^{2}$.
The other side ($A$ is an isometry)
\[
\left|\left|x+y\right|\right|^{2}=\left|\left|x\right|\right|^{2}+\left|\left|y\right|\right|^{2}+2\left\langle x,y\right\rangle 
\]
 we also know that $\left|\left|x+y\right|\right|^{2}=\left|\left|A\left(x+y\right)\right|\right|^{2}=\left|\left|Ax+Ay\right|\right|^{2}=\left|\left|Ax\right|\right|^{2}+\left|\left|Ay\right|\right|^{2}+2\left\langle Ax,Ay\right\rangle $
we know that $\left|\left|Ax\right|\right|^{2}=\left|\left|x\right|\right|^{2},\left|\left|Ay\right|\right|^{2}=\left|\left|y\right|\right|^{2}$.
So 
\begin{align*}
\left|\left|Ax\right|\right|^{2}+\left|\left|Ay\right|\right|^{2}+2\left\langle Ax,Ay\right\rangle  & =\left|\left|x+y\right|\right|^{2}=\left|\left|x\right|\right|^{2}+\left|\left|y\right|\right|^{2}+2\left\langle x,y\right\rangle \\
2\left\langle Ax,Ay\right\rangle  & =2\left\langle x,y\right\rangle \\
\left\langle Ax,Ay\right\rangle  & =\left\langle x,y\right\rangle 
\end{align*}
\end{proof}
We saw: if $A$ is an isometry $\forall x,y$. $\left\langle Ax,Ay\right\rangle =\left\langle x,y\right\rangle $
\[
\left\langle A^{T}Ax,y\right\rangle =\left\langle x,y\right\rangle 
\]
$fix$ x $\forall y$ $\left\langle A^{T}Ax-x,y\right\rangle =0\Rightarrow$$A^{T}A=I$.%
\begin{insight}
because for all $y$ it happens that $A^{T}Ax=x$ which is $\left\langle A^{T}Ax-x,y\right\rangle =\left\langle x-x,y\right\rangle =0$
this wasn\'{t} the case if $A^{T}Ax\ne x$.%
\end{insight}

\end{proposition}


\section{Singular value decomposition}

\begin{insight}
spectral: overall name that speaks on eigenvalues and eigenvectors%
\end{insight}

\begin{thm}[The spectral theorem for real symmetric matrices]
 Every real symmetry matrix $A$ can be expressed as: $A=UDU^{T}$
where $U$ is orthogonal, $D$ is diagonal on $D's$ diagonal are
the eigenvalues of $A$. (Without proof) %
\begin{insight}
Note that $U^{T}AU=D$%
\end{insight}
\end{thm}

let us remember again the operator norm:
\begin{definition}
\begin{insight}
Again (note that you defined it back there\.{)}%
\end{insight}
{} Let there be a real matrix $A$ (not necessarily square), it's operator
norm is $\left|\left|A\right|\right|_{op}=\max\limits _{x}\frac{\left|\left|Ax\right|\right|}{\norm x}=\max\limits _{\left|\left|x\right|\right|=1}\left|\left|Ax\right|\right|$.
\end{definition}

\begin{thm}[SVD -- Singular Value Decomposition]
 Every real matrix $A_{m\times n}$ can be expressed as 
\[
A_{m\times n}=U_{m\times m}D_{m\times n}V_{n\times n}^{T}
\]
 where $U$ and $V$ are orthogonal, and $D$ is ``diagonal''. 
\[
D=\left(\begin{matrix}\sigma_{1} & \dots & 0 & 0\\
0 & \sigma_{2} & \dots & 0\\
\vdots &  & \ddots & \vdots\\
0 & \dots &  & \sigma_{n}\\
0 & 0 & \dots & 0\\
\vdots &  &  & \vdots\\
0 & 0 & \dots & 0
\end{matrix}\right),m\ge n
\]
where $\sigma_{1}\ge\sigma_{2}\ge\dots\ge0$ $'A'$s singular values.
\end{thm}

Note that for $n\ge m$ $D=\left(\begin{matrix}\sigma_{1} & \dots & 0 & 0 & 0 & \dots & 0\\
0 & \sigma_{2} & \dots & 0 & 0 &  & 0\\
\vdots &  & \ddots & \vdots & 0 &  & 0\\
0 & \dots &  & \sigma_{n} & 0 & \dots & 0
\end{matrix}\right),m\ge n$

\begin{proof}
The proof is by induction on $m+n$. %
\begin{insight}
Base: $m=n=1$ thus $A$ is a $1\times1$ matrix. so let us use 
\[
I_{1}AI_{1}
\]
 $A$ is diagonal and it's eigenvalue is it's own value. Step%
\end{insight}
\begin{insight}
We use induction on $m+n$. Main idea: We want $U^{T}AV=D$, we will
do it in stages, not all at once. First stage : We seek orthogonal
matrices $U_{1},V_{1}$ such that 
\[
U_{1}^{T}AV_{1}=\left(\begin{array}{c|ccc}
\text{real number} & 0 & \dots & 0\\
\hline 0\\
\vdots &  & B\\
0
\end{array}\right)
\]
We want to apply induction to repeat this process and then the process
should do it for us.. We take the decomposition of the orthogonal
to those that take the key roles of $SVD$ in $B$. and by induction
we are done. $\tilde{U}B\tilde{V}^{T}=\tilde{D}$. so we take 
\[
\left(\begin{array}{c|ccc}
1 & 0 & \dots & 0\\
\hline 0\\
\vdots &  & \tilde{U}\\
0
\end{array}\right)
\]
 to get the wanted $U_{2}$ and $V_{2}$.%
\end{insight}

We will find $U_{1},V_{1}$ orthogonal s.t. 
\[
U_{1}^{T}AV_{1}=\left(\begin{array}{c|ccc}
\left|\left|A\right|\right|_{op} & 0 & \dots & 0\\
\hline 0\\
\vdots &  & B\\
0
\end{array}\right)
\]
First lets figure out where $\left|\left|A\right|\right|_{op}$ comes
from. $\exists v_{1},u_{1}$ , $\left|\left|v_{1}\right|\right|=\left|\left|u_{1}\right|\right|=1$
s.t. $Av_{1}=\underbrace{\left|\left|A\right|\right|_{op}}_{\sigma_{1}}u_{1}$
(From the definition of the operator norm%
\begin{insight}
Actually Nati didn'\'{t} show this, 
\[
u_{1}=\frac{Ax}{\left|\left|Ax\right|\right|},v_{1}=x
\]
 and $x$ is the vector that gains $\left|\left|A\right|\right|_{op}$
so $\left|\left|v_{1}\right|\right|=1\iff v=\frac{x}{\left|\left|x\right|\right|}$
and we get $u_{1},v_{1}$ as wanted.%
\end{insight}
). $V_{1}=\left(v_{1}|V_{2}\right)$ (we extend $v_{1}$ into an orthogonal
matrix)%
\begin{insight}
(meaning that $\left(v_{1}|V_{2}\right)$ is an orthogonal matrix
where it's first row is $v_{1}$)%
\end{insight}
{} and $U_{1}=\left(u_{1}|U_{2}\right)$. So 
\begin{align*}
U_{1}^{T}AV_{1} & =\left(\begin{array}{c}
u_{1}^{T}\\
\hline U_{2}^{T}
\end{array}\right)A\left(\begin{array}{c|c}
v_{1} & V_{2}\end{array}\right)=\\
 & =\left(\begin{array}{c}
u_{1}^{T}\\
\hline U_{2}^{T}
\end{array}\right)\left(A\begin{array}{c|c}
v_{1} & AV_{2}\end{array}\right)=\\
 & =\left(\begin{array}{c|c}
u_{1}^{T}Av_{1} & u_{1}^{T}AV_{2}\\
\hline U_{2}^{T}Av_{1} & U_{2}^{T}AV_{2}
\end{array}\right)
\end{align*}
 so $u_{1}^{T}Av_{1}=\left|\left|A\right|\right|_{op}$, $U_{2}^{T}Av_{1}=0$
because $U_{2}^{T}\cdot\left(\left|\left|A\right|\right|_{op}u_{1}\right)=0$
because $U_{2}$ is orthogonal %
\begin{insight}
$u_{1}$ is some basis vector from the orthogonal matrix which means
that $\left\langle u_{1},u_{i}\right\rangle =0$%
\end{insight}
. Currently: we found orthogonal matrices $U_{1},V_{1}$ s.t 
\[
U_{1}^{T}AV=\left(\begin{array}{ccc|ccc}
 & \sigma_{1} &  & w_{1}\\
\hline  & 0 & \\
 & \vdots &  & B\\
 & 0 & 
\end{array}\right)
\]
 and we want to show that $w_{1}=0$. The plan: Show that if $w_{1}\ne0$
then we can find a vector that $A$ extends by a vector $>\sigma_{1}$
\begin{insight}
(in contradiction)%
\end{insight}
. Observe that if a matrix is multiplied by an orthogonal matrix either
from left or from the right, the operator norm does not change.
\[
\left(\begin{array}{ccc|ccc}
 & \sigma_{1} &  & w_{1}\\
\hline  & 0 & \\
 & \vdots &  & B\\
 & 0 & 
\end{array}\right)\cdot\left(\begin{array}{c}
\sigma_{1}\\
\hline \\
w_{1}^{^{T}}\\
\\
\end{array}\right)=\left(\begin{array}{c}
\sigma_{1}^{2}+\left|\left|w_{1}\right|\right|^{2}\\
\\
\vdots\\
\\
\end{array}\right)
\]
 This extends by a factor $\ge\sqrt{\sigma_{1}^{2}+\left|\left|w_{1}\right|\right|^{2}}\ge\sigma_{1}$.
What we conclude from that: 
\[
\sigma_{1}=\left|\left|A\right|\right|_{op}=\left|\left|U_{1}^{T}AV_{1}\right|\right|_{op}\ge\frac{\left|\left|\left(\begin{array}{c}
\sigma_{1}^{2}+\left|\left|w_{1}\right|\right|^{2}\\
\\
\vdots\\
\\
\end{array}\right)\right|\right|}{\sqrt{\sigma_{1}^{2}+\left|\left|w_{1}\right|\right|^{2}}}\ge\sqrt{\sigma_{1}^{2}+\left|\left|w_{1}\right|\right|^{2}}\underbrace{>}_{\text{*}}\sigma_{1}
\]
\end{proof}
\begin{itemize}
\item From the assumption that $w_{1}\ne0$ so we got a contradiction.
\end{itemize}
Hence $w_{1}=0$ and we got what we wanted. now we can apply this
idea using induction. Finishing by induction: We found $U_{1},V_{1}$
orthogonal s.t 
\[
U_{1}^{T}AV_{1}=\left(\begin{array}{ccc|ccc}
 & \sigma_{1} &  & 0 & \dots & 0\\
\hline  & 0 & \\
 & \vdots &  &  & B\\
 & 0 & 
\end{array}\right)
\]
 and $B=\tilde{U}\tilde{D}\tilde{V}^{T}$. So 
\begin{align*}
A & =U_{1}\left(\begin{array}{ccc|ccc}
 & 1 &  & 0 & \dots & 0\\
\hline  & 0 & \\
 & \vdots &  &  & \tilde{U}\\
 & 0 & 
\end{array}\right)\left(\begin{array}{ccc|ccc}
 & \sigma_{1} &  & 0 & \dots & 0\\
\hline  & 0 & \\
 & \vdots &  &  & \tilde{D}\\
 & 0 & 
\end{array}\right)\left(\begin{array}{ccc|ccc}
 & 1 &  & 0 & \dots & 0\\
\hline  & 0 & \\
 & \vdots &  &  & \tilde{V}\\
 & 0 & 
\end{array}\right)V_{1}^{T}\\
U_{1}^{T}AV & =\left(\begin{array}{ccc|ccc}
 & 1 &  & 0 & \dots & 0\\
\hline  & 0 & \\
 & \vdots &  &  & \tilde{U}\\
 & 0 & 
\end{array}\right)\left(\begin{array}{ccc|ccc}
 & \sigma_{1} &  & 0 & \dots & 0\\
\hline  & 0 & \\
 & \vdots &  &  & \tilde{D}\\
 & 0 & 
\end{array}\right)\left(\begin{array}{ccc|ccc}
 & 1 &  & 0 & \dots & 0\\
\hline  & 0 & \\
 & \vdots &  &  & \tilde{V}\\
 & 0 & 
\end{array}\right)
\end{align*}
\begin{insight}
That we get in this process 
\[
\left(\begin{array}{ccc|ccc}
 & \sigma_{1} &  & w_{1}\\
\hline  & 0 & \\
 & \vdots &  & B\\
 & 0 & 
\end{array}\right)\to\left(\begin{array}{ccc|c|cc}
 & \sigma_{1} &  & w_{1}\\
\hline  & 0 &  & \sigma_{2}\\
\hline  & \vdots &  &  & C\\
 & 0 &  & 
\end{array}\right)
\]
and $\sigma_{1}\ne\sigma_{2}$ because in the next step we find $\left|\left|B\right|\right|_{op}$
and not $\left|\left|A\right|\right|_{op}$ again. So we get $\left|\left|A\right|\right|_{op},\left|\left|B\right|\right|_{op},\left|\left|C\right|\right|_{op}\dots$%
\end{insight}

Suppose that indeed $A=UDV^{T}$ as in the theorem, how to interpret
the numbers $\sigma_{1}\ge\sigma_{2}\ge\dots\ge0$? 
\begin{align*}
AA^{T} & =UD\underbrace{V^{T}V}_{I}D^{T}U^{T}=\\
 & =UDD^{T}U^{T}
\end{align*}
\[
DD^{T}=\left(\begin{matrix}\sigma_{1}^{2} & \dots & 0 & 0\\
0 & \sigma_{2}^{2} & \dots & 0\\
\vdots &  & \ddots & \vdots\\
0 & \dots &  & \sigma_{n}^{2}
\end{matrix}\right)_{m\times m}
\]
 Based on the spectral theorem for real symmetric matrices (and note
that $AA^{T}$ is r.s%
\begin{insight}
real symmetric%
\end{insight}
) we conclude that $\sigma_{1}^{2},\dots,\sigma_{n}^{2},\underbrace{0,\dots,0}_{m-n}$
are the eigenvalues of the (Positive semi-definite) real matrix $AA^{T}$.

We take another look at the proof of the SVD theorem to make sure
it is understood by everyone.
\begin{thm}[SVD]
 Every real $m\times n$ matrix $A$ can be expressed as $A=UDV^{T}$
where $U_{m\times m}$ is an orthogonal matrix, $V_{n\times n}$ is
orthogonal and $D$ is ``diagonal'' and on the diagonal of $D$
$\sigma_{1}\ge\dots\ge\sigma_{\min\left\{ m,n\right\} }\ge0$
\end{thm}

When we look at $U^{T}AV$ we get a ``diagonal'' matrix.

When we prove this statement we want to find $U_{1}AV_{1}^{T}$ such
that they give us 
\[
U_{1}^{T}AV_{1}=\left(\begin{array}{ccc|ccc}
 & \sigma_{1} &  & 0 & \dots & 0\\
\hline  & 0 & \\
 & \vdots &  &  & B\\
 & 0 & 
\end{array}\right)
\]
 and by induction we can get for $B$ orthogonal matrices that get
as what we need. so how do we do it? first we find$u_{1},v_{1}$ such
that $Av_{1}=\sigma u_{1}$and $\left|\left|v_{1}\right|\right|=1=\left|\left|u_{1}\right|\right|$.
Now if $U_{1}$ and $V_{1}$ are defined as in the proof we get
\[
U_{1}^{T}AV_{1}=\left(\begin{array}{ccc|ccc}
 & \sigma_{1} &  &  & w\\
\hline  & 0 & \\
 & \vdots &  &  & B\\
 & 0 & 
\end{array}\right)
\]
How large can we make $\sigma$? it can be as large as $\left|\left|A\right|\right|_{op}$
- the biggest such $\sigma$. and we get 
\[
U_{1}^{T}AV_{1}=\left(\begin{array}{ccc|ccc}
 & \left|\left|A\right|\right|_{op} &  &  & w\\
\hline  & 0 & \\
 & \vdots &  &  & B\\
 & 0 & 
\end{array}\right)
\]
 our claim is that for $\sigma=\left|\left|A\right|\right|_{op}$
then $w$ is 0. else we get $\left|\left|A\right|\right|_{op}>\sigma$
which is impossible. So we lets choose some specific vector to get
a contradiction 
\[
\left(\begin{array}{c}
\left|\left|A\right|\right|_{op}\\
\hline w^{T}
\end{array}\right)
\]
Which leads to:
\[
\sigma_{1}=\left|\left|A\right|\right|_{op}=\left|\left|U_{1}^{T}AV_{1}\right|\right|_{op}\ge\frac{\left|\left|\left(\begin{array}{c}
\sigma_{1}^{2}+\left|\left|w_{1}\right|\right|^{2}\\
\\
\vdots\\
\\
\end{array}\right)\right|\right|}{\sqrt{\sigma_{1}^{2}+\left|\left|w_{1}\right|\right|^{2}}}\ge\sqrt{\sigma_{1}^{2}+\left|\left|w_{1}\right|\right|^{2}}\underbrace{>}_{\text{*}}\sigma_{1}
\]
 so $w=0$. and we're done.

Question: Why does the sub matrix $B$ has a smaller operator norm?Working
with $B$ is the same as working with matrix $A$ but with a row filled
with zeros, thus it has a smaller or equal operator norm.

Note $A=UDV^{T}\iff A=\sum\limits _{i=1}^{\min\left\{ m,n\right\} }\sigma_{i}u_{i}\otimes v_{i}$
meaning the outer product of $u_{i}$and $v_{i}$ where $u_{i}$ is
the ith column of $U$ and $v_{i}$ is the ith column %
\begin{insight}
column because we have there $V^{T}$%
\end{insight}
of $V$. 
\[
UDV^{T}=\left(\begin{matrix}| &  & |\\
u_{1} & \dots & u_{m}\\
| &  & |
\end{matrix}\right)D\left(\begin{matrix}- & v_{1} & -\\
 & \vdots\\
- & v_{n} & -
\end{matrix}\right)
\]


\subsection{another operator norm}
\begin{xca}
How to evaluate how ``close'' is matrix $B$ to matrix $A$?
\begin{sol}
two answers:
\begin{enumerate}
\item $\left|\left|A-B\right|\right|_{op}$
\item $\left|\left|A-B\right|\right|_{2}=\left|\left|A-b\right|\right|_{\text{Frobenius}}=\sqrt{\sum\left(a_{ij}-b_{ij}\right)^{2}}$
\end{enumerate}
\end{sol}

\end{xca}

$\left|\left|A\right|\right|_{frobenius}=\sqrt{\sum a_{ij}^{2}}$
note that $\left|\left|A\right|\right|_{fr}^{2}=trace\left(AA^{T}\right)$.
because:
\begin{align*}
\left|\left|A\right|\right|_{fr}^{2} & =trace\left(AA^{T}\right)=\sum_{i}\left(AA^{T}\right)_{ii}\\
\end{align*}
 
\[
\left(AA^{T}\right)_{ii}=\left|\left|a_{i}\right|\right|^{2}
\]
 where $a_{i}$ is the ith row of $A$ so 
\[
\left|\left|a_{i}\right|\right|^{2}=\sum_{j}a_{ij}^{2}
\]
 %
\begin{insight}
\[
trace\left(AA^{T}\right)=\sum_{i,j}a_{ij}^{2}
\]
\end{insight}

\begin{claim}
Let $A$ be a real matrix, $U$ an orthogonal matrix, note that 
\[
\left|\left|A\right|\right|_{fr}=\left|\left|AU\right|\right|_{fr}
\]
\begin{proof}
$ $ 
\[
\left|\left|AU\right|\right|_{fr}^{2}=tr\left(\left(AU\right)\left(AU\right)^{T}\right)=tr\left(AUU^{T}A^{T}\right)=tr\left(AIA^{T}\right)=tr\left(AA^{T}\right)=\left|\left|A\right|\right|_{fr}^{2}
\]
\end{proof}
\end{claim}

\begin{thm}
Let $A_{m\times n}$ be a real matrix and let $k<m,n$. For every
$B_{m\times n}$ matrix of rank $\le k$ there holds $\left|\left|A-B\right|\right|_{op}\ge\sigma_{k+1}$
where $\sigma_{1}\ge\dots\ge0$ are the singular values of $A$. Equality
holds for$B=UD^{\left(k\right)}V^{T}$, where $A=UDV^{T}$ and 
\[
D^{\left(k\right)}=\left(\begin{array}{cccccc}
\sigma_{1} & 0 & \dots &  & \dots & 0\\
0 & \ddots &  &  &  & \vdots\\
\vdots &  & \sigma_{k}\\
 &  &  & 0\\
 &  &  &  & \ddots & \vdots\\
\vdots &  &  &  &  & 0\\
0 &  &  &  &  & 0
\end{array}\right)
\]
 
\[
B=\sum_{i=1}^{k}\sigma_{i}u_{i}\otimes v_{i}
\]
(We can\'{t} do better then than $\left|\left|A-B\right|\right|_{op}\ge\sigma_{k+1}$)
\end{thm}

Note this: $A-B_{k}=U\left(\begin{array}{ccccc}
0\\
 & \ddots\\
 &  & 0\\
 &  &  & \sigma_{k+1}\\
 &  &  &  & \ddots
\end{array}\right)V^{T}\Rightarrow$$\left|\left|A-B\right|\right|_{op}=\sigma_{k+1}$
\begin{proof}
\begin{insight}
all Lecture page 31%
\end{insight}
{} What do we need to show? for $C$ matrix such that $\text{rank}\left(C\right)\le k\Rightarrow\left|\left|A-C\right|\right|_{op}\ge\sigma_{k+1}$.
\begin{insight}
to show this we need to show some vector $z$ such that $\left|\left|\left(A-C\right)z\right|\right|\ge\sigma_{k+1}$%
\end{insight}
Lets start: $\exists z,\left|\left|\left(A-C\right)z\right|\right|_{2}\ge\sigma_{k+1}$
where $\left|\left|z\right|\right|_{2}=1$. How can we find such a
vector? Plan: We will select $z$ in such a way that $Cz=0$ and then
$\left|\left|\left(A-C\right)z\right|\right|=\left|\left|Az-Cz\right|\right|=\left|\left|Az\right|\right|$
- lets do just that: Since $\text{rank}\left(C\right)\le k$ the kernel
of $C$ intersects (not in $\bar{0}$) every linear subspace $W$
of $\dim W\ge k+1$. Let $W=span\left\{ v_{1}\dots v_{k+1}\right\} $
(the columns of $V$).
\[
z=\sum_{i=1}^{k+1}\alpha_{i}v_{i};Cz=0;\left|\left|z\right|\right|_{2}=1
\]
 
\begin{align*}
\left(A-C\right)z & =Az-Cz=Az=\sum_{i}\sigma_{i}\left(u_{i}\otimes v_{i}\right)z=\\
 & =\sum\sigma_{i}u_{i}\left\langle v_{i},z\right\rangle =\sum^{m}\sigma_{i}u_{i}\left\langle v_{i},z\right\rangle \\
\left(1\right) & =\sum_{i=1}^{k+1}\sigma_{i}u_{i}\left\langle v_{i},z\right\rangle =\sum_{i=1}^{k+1}\sigma_{i}\left\langle v_{i},z\right\rangle u_{i}=
\end{align*}
\begin{enumerate}
\item everything beyond $k+1$ equals zero. because $v_{1}\dots v_{m}$are
orthonormal then $\left\langle v_{i},z\right\rangle =\left\langle v_{i},\sum_{j}^{k}\alpha_{j}v_{j}\right\rangle $
equals 0 for every $j\ne i$.
\end{enumerate}
Because the $u_{i}$'s are orthogonal:
\begin{align*}
\left|\left|\left(A-C\right)z\right|\right|^{2} & =\sum_{i=1}^{k+1}\sigma^{2}\left\langle v_{i},z\right\rangle ^{2}\ge\sigma_{k+1}^{2}\underbrace{\sum\limits _{i=1}^{k+1}\left\langle v_{i},z\right\rangle ^{2}}_{(*)}\\
 & =\sigma_{k+1}^{2}\cdot1
\end{align*}

\begin{itemize}
\item because it's an orthogonal basis.%
\begin{insight}
To my understanding because it\'{s} an orthonormal basis%
\end{insight}
\end{itemize}
\end{proof}
%
\begin{definition}
A Real symmetric matrix is said to be positive semi-definite (PSD)
if all it's eigenvalues are non negative.
\end{definition}

\begin{definition}
A Real symmetric matrix is said to be positive definite (PD) if all
it's eigenvalues are positive.
\end{definition}

\begin{thm}
The following are equivalent (TFAE) regarding a real symmetric matrix
$A$
\begin{enumerate}
\item $A$ is PSD.
\item $A=MM^{T}$ for some real $M$.
\item $\forall x,xAx^{T}\ge0$
\end{enumerate}
\end{thm}

\begin{proof}
$ $
\end{proof}
\begin{itemize}
\item[$1\Rightarrow2$:]  recall the spectral theorem for real symmetric matrices: it is possible
to express $A=U\Lambda U^{T}$ where $U$ is orthogonal $\Lambda$
is diagonal, the diagonal terms in $\Lambda$ are $A$'s eigenvalues
and the row of $U$ are the corresponding eigenvectors. Let $M=U\sqrt{\Lambda}$
($\sqrt{\Lambda}$ = The diagonal matrix whose $i$ths entry is $\sqrt{\lambda_{i}}$.)
$MM^{T}=U\sqrt{\Lambda}\sqrt{\Lambda}U^{T}=U\Lambda U^{T}=A$. As
wanted.
\item[$2\Rightarrow3$:]  Since $A=MM^{T}$ $xAx^{T}=xMM^{T}x=\left|\left|xM\right|\right|_{2}^{2}\ge0$,
As wanted.
\item[$3\Rightarrow1$:]  Let $v$ be an eigenvector with eigenvalue $\lambda$. We know that
$vAv^{T}\ge0$ but on the other hand 
\[
0\le vAv^{T}=\lambda\left\langle v,v\right\rangle =\lambda\left|\left|v\right|\right|_{2}^{2}
\]
\begin{insight}
Notice that the rows are revresed: (full explanation for the equation)
\[
0\le vAv^{T}=v\left(\lambda v^{T}\right)=\lambda\left\langle v,v\right\rangle =\lambda\left|\left|v\right|\right|_{2}^{2}
\]
\end{insight}
{} so $\lambda\ge0$ because $\left|\left|v\right|\right|_{2}^{2}\ge0$.
\end{itemize}
\begin{corollary}
The collection of all $n\times n$ PSD matrices form a \textbf{cone}
$PSD_{n}$. Namely if $A,B\in PSD_{n}\Rightarrow A+B\in PSD_{n}$
and if $t>0\Rightarrow tA\in PSD_{n}$.%
\begin{insight}
We can use the 3 property we proved to show that.%
\end{insight}
\end{corollary}


\subsection{The variational definition of eigenvalues of real symmetric matrices}
\begin{thm}[Rayleigh-Ritz]
 Let $A$ be real symmetric matrix and let $\lambda_{1}\ge\lambda_{2}\ge\dots\ge\lambda_{n}$
be it's eigenvalues. Then 
\[
\lambda_{1}=\max\frac{xAx^{T}}{\left|\left|x\right|\right|_{2}^{2}}=\max\limits _{\left|\left|x\right|\right|_{2}=1}\left(xAx^{T}\right)
\]
$\lambda_{2}=\max\limits _{x\perp x_{1}}\frac{xAx^{T}}{\left|\left|x\right|\right|_{2}^{2}}$,
$\lambda_{3}=\max\limits _{x\perp x_{1}\perp x_{2}}\frac{xAx^{T}}{\left|\left|x\right|\right|_{2}^{2}}$,
and so on. 
\end{thm}

To prove that we will use some analysis:
\begin{remark}
Recall that $\left(\frac{f}{g}\right)^{'}=0\iff fg'=f'g$ we'll have
$f=\left(xAx^{T}\right)$ and $g=\left|\left|x\right|\right|^{2}$
and we want to get the maximum of that.
\end{remark}

\begin{remark}
$\left(xA\right)_{\nu}=\sum\limits _{s}a_{s\nu}x_{s}$ and $xAx^{T}=\sum_{\alpha,\beta}x_{\alpha}a_{\alpha,\beta}x_{\beta}$.
\end{remark}

We will show that if $x$ maximizes $\frac{xAx^{T}}{\left|\left|x\right|\right|_{2}^{2}}$
then $x$ is an eigenvector. First lets derive $\left|\left|x\right|\right|_{2}^{2}$:
$\frac{\partial}{\partial x_{i}}\left|\left|x\right|\right|_{2}^{2}=2x_{i}$.
Now from the previous remark:
\[
\frac{\partial}{\partial x_{i}}xAx^{T}=\frac{\partial}{\partial x_{i}}\left(\sum\limits _{\alpha,\beta}a_{\alpha,\beta}x_{\alpha}x_{\beta}\right)=2\left(Ax\right)_{i}
\]
 twice is because the row can either come from $a_{ij}x_{i}x_{j}$
or from $a_{si}x_{s}x_{i}$. for $x$ to maximize the Rayleigh quotient
$\frac{xAx^{T}}{\left|\left|x\right|\right|^{2}}$ a necessary condition
is that $\forall i$ $\underbrace{\left(Ax\right)_{i}}_{f'}\cdot\underbrace{\left|\left|x\right|\right|^{2}}_{g}=\underbrace{\left(xAx^{t}\right)}_{f}\underbrace{x_{i}}_{g'}$
notice that $\left|\left|x\right|\right|^{2},\left(xAx^{T}\right)$
are numbers, and 
\[
\underbrace{\left(Ax\right)_{i}}_{f'}\cdot\underbrace{\left|\left|x\right|\right|^{2}}_{g}=\underbrace{\left(xAx^{t}\right)}_{f}\underbrace{x_{i}}_{g'}\iff\left(Ax\right)\left|\left|x\right|\right|^{2}=\left(xAx^{T}\right)x
\]
so we've seen that $\left(Ax\right)\underbrace{\left|\left|x\right|\right|^{2}}_{number}=\underbrace{\left(xAx^{T}\right)}_{number}x$.
We've seen that if $x$ maximizes $\frac{xAx^{T}}{\left|\left|x\right|\right|_{2}^{2}}$
then $x$ is an eigenvector.

\begin{proof}[proof of Rayleigh - Ritz:]
 (We prove only for $\lambda_{1}$, the next steps ($\lambda_{2}\dots$)
need to use Lagrange multipliers. %
\begin{insight}
The concept is that he shows here that if x maxsimizes this (f is
the nominator and $g$ is the denominator of what we want) so if $x$
maximises it then it is an eigenvalue!! So if we maximise this thing
we should find the biggest lambda! (maximal x is always eigenvector
so we use maximal eigenvalue to get the maximum of $\frac{xAx^{T}}{\left|\left|x\right|\right|_{2}^{2}}$)%
\end{insight}
) Using the spectral theorem: $A=V\Lambda V^{T}\Rightarrow$$xAx^{T}=xV\Lambda V^{T}x^{T}$
We think of it as $y\Lambda y^{T}=\sum\lambda_{i}y_{i}^{2}$ where
$y=xV^{T}$ So 
\[
=\underbrace{\sum\lambda_{i}\left(Vx\right)_{i}^{2}}_{\text{positive items}}\le\left(\max\lambda_{i}\right)\cdot\sum\left(Vx\right)_{i}^{2}=\lambda_{1}\sum\left(Vx\right)_{i}^{2}=\lambda_{1}\left|\left|Vx\right|\right|^{2}=\lambda_{1}\left|\left|x\right|\right|^{2}
\]
 if $Ax=\lambda_{1}x$ then $\frac{xAx^{T}}{\left|\left|x\right|\right|^{2}}=\lambda_{1}$%
\begin{insight}
for $\lambda_{2}$: The same thing but if $x$ is orthogonal to the
previous $x$ then the equations zeros out.%
\end{insight}
\end{proof}
%
\begin{thm}[Rayliegh-Ritz 2]
 Let $A$ be a real symmetric matrix with eigenvalues $\lambda_{1}\ge\lambda_{2}\ge\dots\ge\lambda_{n}$
. Then$\lambda_{n}=\min\frac{xAx^{T}}{\left|\left|x\right|\right|^{2}};$$\lambda_{n-1}=\min\limits _{x\perp x_{n}}\frac{xAx^{T}}{\left|\left|x\right|\right|^{2}}$.
(With no proof)
\end{thm}

How do we derive the second theorem from the first? we look at the
minus of the same thing.
\begin{thm}[Courant-Fischer]
 $A_{n\times n}$ real symmetric, with eigenvalues $\lambda_{1}\ge\lambda_{2}\ge\dots\ge\lambda_{n}$.
Then $\forall i$ 
\[
\lambda_{i+1}=\min_{\begin{array}{c}
F\\
\dim F=i
\end{array}}\max_{x\perp F}\left\{ \frac{xAx^{T}}{\left|\left|x\right|\right|^{2}}\right\} 
\]
 and $\max\limits _{\dim F=j}\min\limits _{x\perp F}\frac{xAx^{T}}{\left|\left|x\right|\right|^{2}}=\lambda_{n-j}$.
\begin{insight}
Why is there an optimum for all these subspaces? (There's infinite
amount of subspaces) Because of compactness.%
\end{insight}
\end{thm}

\begin{thm}[The Interlacing theorem]
Let $A$ be a real symmetric matrix with eigenvalues $\lambda_{1}\ge\dots\ge\lambda_{n}$.
Let $B$ be the matrix that results by erasing the i-th row and the
i-th column of $A$. %
\begin{insight}
B is real symmetric too. and has $n-1$ eigenvalues.%
\end{insight}
{} Let $\mu_{1}\ge\dots\ge\mu_{n-1}$ be the eigenvalues of $B$. So
$\lambda_{1}\ge\mu_{1}\ge\lambda_{2}\ge\mu_{2}\ge\dots\ge\lambda_{n-1}\ge\mu_{n-1}\ge\lambda_{n}$.
\begin{insight}
This is derived from Rayleigh -- Ritz. We need to connect $A$ to
$B$. so the way to think about this is that we can pretendt that
$B$ was never born and we keep working with $A$. we insist that
the $i'th$ coordinate of $x$ is 0. Doing so we get what we get $B$.%
\end{insight}
\end{thm}


\section{Perron -- Frobenius Theorems}
\begin{thm}[Perron]
 Let $A$ be a square positive matrix %
\begin{insight}
positive means its real.%
\end{insight}
(i.e $\forall i,j$ $a_{i,j}>0$) Then $A$ has a positive eigenvalue
$\rho>0$ s.t. for every other eigenvalue $\mu$ (in general complex)
$\rho>\left|\mu\right|$. The eigenvector corresponding to $\rho$
is positive, no other eigenvector is positive.
\end{thm}

The proof's stages:
\begin{enumerate}
\item $\rho>0$ corresponding eigenvector $>0$.
\item $\forall\mu\ne\rho,\,\,\rho>\left|\mu\right|$.
\item $\rho$ is a simple eigenvalue. (Has an algebraic multiplicity 1 --
Note that we don't actually show algebraic multiplicity 1) 
\item no other positive eigenvector (We don't prove this)
\end{enumerate}
\begin{insight}
we want a theorem for non negative matrices, rather then strictly
positive matrix. to show that $A^{k}$ is positive we switch every
non zero number to 1. it still holds that if this new matrix is positive
so is $A^{k}$. now that $A$ is zeros or 1s we know that $A$ is
a the adjacency matrix of a directed graph. $Au=\lambda u$ then $A^{k}u=\lambda^{k}u$.
If we knew persons theorem and we know that $A^{k}>0$ we can derive
the same conclusions for Perron's Theorem to a non negative matrix.
$A^{k}\ge0$.%
\end{insight}

\begin{proof}[proof of Perron's Theorem]
\begin{insight}
We want to think about it as an optimization problem.%
\end{insight}
Let us observe: if $v\ge0,v\ne0\Rightarrow Av>0$. Let $S\subseteq\mathbb{R}^{n}$
be defined as: 
\[
S=\left\{ u\in\R^{n}:\left|\left|u\right|\right|_{2}=1,u\ge0\right\} 
\]
\begin{insight}
this is a compact set%
\end{insight}

We define $\forall u\in S$, 
\[
L\left(u\right):=\max\left\{ b>0:\begin{array}{c}
\left(Au\right)_{i}\ge bu_{i},\\
\forall i\,\,u_{i}\ne0
\end{array}\right\} =\min_{\begin{array}{c}
i\\
u_{i}\ne0
\end{array}}\frac{\left(Au\right)_{i}}{u_{i}}
\]
 %
\begin{insight}
The coordinate that holds be to become bigger is the $i$ that makes
this$\frac{\left(Au\right)_{i}}{u_{i}}$ the minimum%
\end{insight}

It can be shown that $L\left(u\right)$ is compact%
\begin{insight}
The edge of a set is a point that contains around it points in the
set and points outside of the set.

A close set is a set such it contains it's edge.

Compact set:
\begin{enumerate}
\item Let there be a metric space $X$. A compact set $Y\subseteq X$ is
called compact if for any series $y_{n}\in Y$ There is a converging
sub series of $y_{n}$ such that it converge into a point in $Y$.
\item Every closed set $Z\subseteq Y$ is compact.
\item a compact set is bound.
\item The image of a compact set by a function on that set is also compact.
\item Any non empty compact set contains a maximum and a minimum.
\item A close and bound set in $\R^{n}$ is compact.
\end{enumerate}
\end{insight}
, hence $\rho:=\max\limits _{u\in S}L\left(u\right)$. Let $v\in S$
be a vector realizing $Av\ge\rho v$%
\begin{insight}
from the L$\left(u\right)$ statement%
\end{insight}
. %
\begin{insight}
we want to show that this is an eigenvalue. but we saw that we can
get a bigger value by contradiction.%
\end{insight}
We want to show $Av=\rho v$, if it is -- were done, else we assume
by contradiction that $Av-\rho v>0$, it follows that %
\begin{insight}
non negative and non zero vector multiplied by a positive matrix is
positive.%
\end{insight}
$A\left(Av-\rho v\right)>0$. Because that is a positive vector $\exists\epsilon>0$
s.t. $A\left(Av-\rho v\right)>\epsilon Av$, Let $c>0$ s.t. $cAv$
is a unit vector, actually $cAv\in S$ so%
\begin{insight}
from this: $Av-\rho v>0$ we expand the term:
\begin{align*}
AAv-Apv & >\epsilon Av\\
cAAv-Acpv & >\epsilon cAv\\
AcAv & >\epsilon cAv+Acpv\\
A\left(cAv\right) & >\left(p+\epsilon\right)cAv
\end{align*}
 Notice that $cAv$ is a vector in $S$ for which we defined $\rho=\max\limits _{u\in S}L\left(u\right)$
which will lead us to a contradiction. HELL YES!%
\end{insight}
\[
A\left(cAv\right)>\left(\rho+\epsilon\right)cAv
\]
{} This is contrary to the maximality of $\rho$.
\begin{claim}
If $\mu\ne\rho$ is an eigenvalue of $A$ then $\rho>\left|\mu\right|$. 
\begin{proof}
Let $Ay=\mu y\iff\forall i\,\,\underbrace{\left(Ay\right)_{i}}_{=\sum\limits _{j}a_{i,j}y_{j}}=\mu y_{i}$
So $\left|\sum\limits _{j}a_{ij}y_{j}\right|=\left|\mu\right|\left|y_{i}\right|$
then from triangle inequality 
\[
\left|\mu\right|\cdot\left|y_{i}\right|=\left|\sum a_{ij}y_{j}\right|\le\sum\left|a_{i,j}\right|\left|y_{j}\right|=\sum a_{i,j}\left|y_{j}\right|
\]
. Denote $z_{i}=\left|y_{i}\right|$ so $\sum a_{i,j}z_{i}\ge\left|\mu\right|\cdot z_{i}$
meaning:
\[
Az\ge\left|\mu\right|z
\]
 From this it follows that $\rho\ge\left|\mu\right|$. %
\begin{insight}
$\rho$ is defined as the largest number such that $Av\ge\rho v$
so $\mu$ cannot be bigger then $\rho$.%
\end{insight}
\begin{insight}
z should be normalized and then $z\in S$ and as a result $\rho$
is the maximal eigenvalue.%
\end{insight}
\end{proof}
\end{claim}

\begin{claim}
$\rho$ has a geometric multiplicity i.e. the linear space of eigenvectors
with eigenvalue $\rho$ is $1-dimensional$
\begin{proof}
suppose that $Au=\rho u$. We can write $u=\phi+i\psi$ and each part
will be an eigenvector separately. This means we can just assume that
$u$ is a real vector. Generally $u$ may be complex, but then the
vectors $Re\left(u\right),Im\left(u\right)$ are also eigenvectors
with eigenvalue $\rho$. WLOG $u$ is real.%
\begin{insight}
We should apply this for both vectors - to the imaginary vector too
and we get the same contradiction.%
\end{insight}

Let us assume by contradiction that $Av=\rho v,Au=\rho u$. $\exists\epsilon$
(might be negative too) s.t. $v+\epsilon u\ge0$ with at least one
zero coordinate. because $v,u$ are eigenvectors of $\rho$: $A\left(v+\epsilon u\right)=\rho\left(v+\epsilon u\right)$.
(This is a contradiction because $A\left(v+\epsilon u\right)$ can't
have an entry that is zero because a positive matrix multiplied by
a non negative vector creates a positive vector -- can't have a cell
with zero in it like the vector $\rho\left(v+\epsilon u\right)$ and
from this we conclude that there can't be another vector $u$ such
that $Au=\rho u$.)

(To my understanding we did not show in this proof that the algebraic
multiplicity is 1)%
\begin{insight}
There cant be an entry with zero and thus the equation is a contradiction
- not sure whats going on. not sure what happened here%
\end{insight}

\end{proof}
\end{claim}

 Additional item in Perron\'{s} Theorem: $A$ positive square matrix
with $\rho$ as before...
\[
\limn{\frac{A}{\rho}}^{n}=\frac{x_{r}\otimes x_{\ell}}{\left\langle x_{r},x_{\ell}\right\rangle }
\]
 ($\otimes$ means outer-product) Where $x_{r},x_{\ell}$, are the
right and left eigenvectors of $A$ corresponding to $\rho$.%
\begin{insight}
We divide each eigenvalue by $\rho^{n}$%
\end{insight}

\end{proof}
%
To pass from Perron to Frobenius: if $A^{^{k}}>0$ for some $k\in\N$
then $A$ satisfies the conclusions of Persons theorem.

Questions: Given $A\ge0$ when does $\exists k$ s.t. $A^{k}>0$ ?
$\iff$ Given a directed graph $G$ under what conditions does $\exists k\in\N$
s.t. $\forall x,y\in V$ $\exists$ directed $x\to y$ of length $k$
.
\begin{thm}
(\textcolor{red}{This wasn't stated in class explicitly}): (Perron
- Frobenius) If $A$ is non negative and there exists a $k\in\N$
such that $A^{k}>0$ then Perron's statement holds for this matrix
too.
\end{thm}

\begin{definition}
We say that directed graph $G=\left(V,E\right)$ has period $p$ if
it is possible to split $V=V_{1}\cup V_{2}\dots\cup V_{p}$ s.t if
$x\to y\in E$ iff $\exists i$ such that $x\in V_{i},y\in V_{i+1\mod p}$.
\end{definition}

\begin{claim}
This is the case provided
\begin{enumerate}
\item $G$ is strongly connected.
\item Aperiodic. %
\begin{insight}
If it's periodic we know that $k$ moves us to some group.%
\end{insight}
\end{enumerate}
\end{claim}


\section{finite Markov chains}

This is a system ($M$) that can reside in any one of $n$ states.
It moves at every time step from state to state, and if it in presently
resides in state $i$ then with probability $p_{i,j}$ it will be
in state $j$ at the next time unit. $P=\left(p_{i,j}\right)$ is
called the transition matrix of the chain. $P$ is a stochastic matrix,
i.e. $\forall i,j$ $p_{i,j}\ge0$ each row sums to one ($\forall i\,\,\sum_{j}p_{ij}=1$).
\begin{definition}
We say that Markov chain is ergodic if the directed graph is strongly
connected and aperiodic.
\end{definition}


\subsubsection{informal chatter}

$P_{i,j}^{t}=$ probability that starting from state $i$ we will
be at state $j$ after $t$ units of time.

Let us restrict ourselves to deal with ergodic Markov chains.

Since the underlying directed graph of $M$ (markov chain) is strongly
connected and aperiodic the Perron -- Frobenius theorem applies.
(This matrix contains one positive eigenvector, other eigenvalues
are strictly smaller then the corresponding eigenvalue) $P\cdot\one=\one\Rightarrow\one$
is an eigenvector corresponding to eigenvalue $1$. By $P.F$ (Perron
-- Frobenius) $\one$ and the eigenvalue $1$ are the eigenvector
and eigenvalue pair given by the theorem $\Rightarrow$ all other
eigenvalues $\lambda$ of $p$ satisfies $\left|\lambda\right|<1$.

Let $u$ be the probability vector(where $u_{i}$ = is the probability
that the system starts at state $i$, and$u\ge0$,$\sum u_{i}=1$):
so 
\[
uP\one=u\one=\left\langle u,\one\right\rangle =1
\]
$\left(uP\right)_{\alpha}=\sum_{i}u_{i}p_{i\alpha}=$probability that
the chain is in state $\alpha$ at the next time unit.

$\left(uP^{2}\right)_{\beta}=\left(\left(uP\right)\cdot P\right)_{\beta}=$
prob that the chain will be in state $\beta$ in two units from now.

What do we want to know using Markov chains?
\begin{enumerate}
\item what is the long-run behavior of the chain.
\item How soon does it get there.
\end{enumerate}
%
From $P.F$ we know that there is a positive \uline{left} eigenvector
$\pi$ corresponding to eigenvalue $1$, $\pi P=\pi$. Let us normalize
$\pi$ so that $\left\langle \pi,\one\right\rangle =1$.
\begin{thm}[from Perron -- frobenious]
 $\limn{\left(\frac{A}{p}\right)^{n}}=\frac{x_{r}\otimes x_{\ell}}{\left\langle x_{r},x_{\ell}\right\rangle }$
where $\rho$ is the eigenvalue from perron--frobenious.$x_{r},x_{\ell}$
are the right and left eigenvectors of $\rho$.
\end{thm}

\begin{proof}
in the case of Markov Chains $A=P,\rho=1$. $x_{r}=\one,x_{\ell}=\pi$$\Rightarrow$
$\limn{\frac{P^{n}}{\rho^{n}}}=\one\otimes\pi=\begin{pmatrix}- & \pi & -\\
- & \pi & -\\
 & \vdots\\
- & \pi & -
\end{pmatrix}$. If $\limn{P^{n}}$ exists let us denote it by $Q$, thus $QP=Q$
$\Rightarrow$ Every row of $Q$ is a left eigenvector of $P$ with
eigenvalue $1$. %
\begin{insight}
(This gives us the matrix above)%
\end{insight}
. Why does the limit exist? The set of all $m\times n$ stochastic
matrices is compact %
\begin{insight}
This a simplex closed and bounded%
\end{insight}
$\Rightarrow$ every infinite sequence of stochastic matrices has
a convergent sub-sequence. As we saw only $Q$ can be a limit point
of matrices in the sequence. %
\begin{insight}
All stochastic matrices is a compact set and we have a sequence so
it has a convergent sequence. Every stochastic matrix will converge
and as a result we know that $Q$ is the only limit? something like
that.%
\end{insight}
\end{proof}
%

\subsection{Random Walks on Graphs}

$G=\left(V,E\right)$ an undirected graph. $States=V$. Transitions
matrix defined is normalized in row $i$ by $\frac{1}{\deg\left(v_{i}\right)}$.
\begin{proposition}
If $G$ is a connected non-bipartite undirected graph, then the limit
distribution of the random walk on $G$ is $\left(\cdots\frac{d\left(x\right)}{2\left|E\right|}\cdots\right)=\pi$.
$\pi P\overset{}{=}\pi$.
\end{proposition}


\subsubsection{Expander graphs}

Expander graphs walks on them, etc.

(informal definition) These are graphs with no bottlenecks. $G=\left(V,E\right)$
d-regular graph. $h\left(G\right):=\min\limits _{\begin{array}{c}
\emptyset\ne S\subseteq V\\
\left|S\right|\le\frac{\left|V\right|}{2}
\end{array}}\frac{e\left(S,\bar{S}\right)}{\left|S\right|}$ %
\begin{insight}
i think the function $e$ counts the number of edges between two sets
of nodes.

Why do we normalize by the size of $S$?%
\end{insight}

We want $d-regular$ graph with a large edge expansion.
\begin{definition}
(TAKEN FROM A PREVIOUS YEAR) Let $\delta\ge0$, a graph $G=\left(V,E\right)$
is $\delta$-edge expanding if for every $S$ of size at most $\frac{\left|V\right|}{2}$
it holds that $e\left(S,\overline{S}\right)\ge\delta\cdot\left|S\right|$.
\end{definition}

\begin{definition}
The edge expansion of a graph $h\left(G\right)$ is the minimal expansion
of a vertices set. I.e. $h\left(G\right)=\min\limits _{\begin{array}{c}
\emptyset\ne S\subset V\\
\left|S\right|\le\frac{\left|V\right|}{2}
\end{array}}\frac{e\left(S,\overline{S}\right)}{\left|S\right|}$
\end{definition}

\uline{Question:} Given $d\ge3$ does there exists $\delta>0$
and infinitely many $d-reg$ graphs $G=G_{n}$ %
\begin{insight}
familiy of $G_{n}$ for $n=d+1\to\infty$%
\end{insight}
with $h\left(G_{n}\right)\ge\delta$?

\begin{insight}
G=Gn - a family for the rest of these claims%
\end{insight}

\begin{thm}
\begin{insight}
\end{insight}
Yes, actually \uline{most} $d-reg$ graphs %
\begin{insight}
family%
\end{insight}
{} have this property.
\end{thm}

For your personal information: To decide whether $h\left(G\right)>\alpha$
is a $co-NP-HARD$ problem.
\begin{thm}
Let $G$ %
\begin{insight}
take a graph - not a family%
\end{insight}
{} be a $d-regular$ graph and let $d=\lambda_{1}\ge\lambda_{2}\ge\dots\ge\lambda_{n}$
be the eigenvalue of adjacency matrix. so %
\begin{insight}
$\Lambda=\max\lambda_{2},-\lambda_{n}$. This is what he wrote then
erased from the board.%
\end{insight}
{} %
\begin{insight}
Note that there must be a negative eigenvalues because the trace of
a graph's matrix is $0$ (on the diagonal is zeros)%
\end{insight}
{} $\frac{d-\lambda_{2}}{2}\le h\left(G\right)\le\sqrt{d^{2}-\lambda_{2}^{2}}$.
\end{thm}

Constructing expander graphs is qualitatively equivalent to constructing
$d-regular$ graphs with $\lambda_{2}$ bounded away from $d$. %
\begin{insight}
building this kind of graph familiy such that $h\left(G_{n}\right)$
is far from zero. equivalent to build a familty of $d-reg$ graphs
where $\lambda_{2}$ is bounded away from $d$. where bounded away
means that there is some small area around $d$ where lambda cannot
reach.%
\end{insight}

\begin{definition}
If $\sigma,\tau$ are two prob distributions on a finite set $S$,
we define their distance 
\[
d_{TV}\left(\sigma,\tau\right)=\sum_{x\in S}\left|\sigma\left(x\right)-\tau\left(x\right)\right|
\]

Please Note that in the recitation we defined it as follows:
\[
d_{TV}\left(\sigma,\tau\right)=\frac{1}{2}\sum_{x\in S}\left|\sigma\left(x\right)-\tau\left(x\right)\right|=\frac{1}{2}\norm{\sigma-\tau}_{1}
\]
\end{definition}

Let $G$ be a regular graph, Let $\sigma_{0}=\left(1,0,\dots,0\right)$,
$\sigma_{t}=$ distribution of the random walker at time $t$.

$Q:$ given $\epsilon>0$, for which time $t$ does it hold that $d_{TV}\left(\sigma_{t},uniform\right)<\epsilon$?
Let $\mathcal{G}=\left(G_{n}\right)$ be a family of $d-regular$
graphs, then the following are equivalent
\begin{enumerate}
\item $h\left(G_{n}\right)$ is bounded away from zero
\item $d-\lambda_{2}$ is bounded away from zero.
\item mixing time is $O\left(\log n\right)$
\end{enumerate}
Let $A$ be a positive matrix $\rho=$ spectral radius of $A$. $\limn{\frac{A}{\rho}}^{n}=\frac{x_{R}\otimes x_{L}}{\left\langle X_{R},X_{L}\right\rangle }$
where $x_{R},x_{L}$ are the right and left eigenvectors of $\rho$.
We will prove the case of $A$ a stochastic matrix.
\begin{proof}
$ $
\begin{lemma}
Let $G$ be a strongly connected aperiodic directed graph and let
$A$ be it's adjacency matrix. Then $\exists k$ s.t $A^{k}>0$. (Without
proof).
\end{lemma}

We will prove the case of a positive stochastic matrix $\then$ same
for the stochastic ergodic case.

Let $A$ be a positive stochastic matrix and let $\alpha>0$ be the
smallest entry in $A$. Let $v$ be a probability vector $v\ge0$,
$\left\langle v,\one\right\rangle =1$. Let $M,m$ be the largest
and smallest coordinates of $v$. Let $M',m'$ be the max/min of $Av$.
\begin{align*}
M' & \le\alpha m+\left(1-\alpha\right)M\\
m' & \ge\alpha M+\left(1-\alpha\right)m
\end{align*}
\begin{insight}
We are multiplying a row in $A$ by a vector. We engineer a vector
that sums to one and gets the biggest number. what do we do? alpha
will do the least damage when it is multiplied by $m$, now after
we do $\alpha m$ we add $\left(1-\alpha\right)M$ and we ignore all
other values, %
\end{insight}
{} so 
\[
M'-m'\le\left(1-2\alpha\right)\left(M-m\right)
\]
and $A^{n}v=A^{n-1}\left(Av\right)=A^{n-2}A\left(Av\right)$ $\then$
$M'^{\left(n\right)}-m'^{\left(n\right)}\le\left(1-2\alpha\right)^{n}\left(M-m\right)$
\[
\limn{M'^{\left(n\right)}-m'^{\left(n\right)}}=0
\]
 then 
\[
\limn{A^{n}v}=c\cdot\one
\]
 for some $c>0$. $v=e_{i}\then$ The col of $\limn{A^{n}}$ is a
constant vector. Therefore $\limn{A^{n}}=\begin{pmatrix}- & u & -\\
- & u & -\\
 & \vdots\\
- & u & -
\end{pmatrix}$ and we've seen that if there such limit it equals to $=\begin{pmatrix}- & \pi & -\\
- & \pi & -\\
 & \vdots\\
- & \pi & -
\end{pmatrix}$ %
\begin{insight}
(from 106)%
\end{insight}
. $c=\left\langle u,v\right\rangle =\left\langle \pi,v\right\rangle $
for each $e_{i}$, $c$ is different.
\end{proof}
%
Edge extension of $G$: $h\left(G\right):=\min\limits _{\begin{array}{c}
\emptyset\ne S\subseteq V\\
\left|S\right|\le\frac{\left|V\right|}{2}
\end{array}}\frac{e\left(S,\overline{S}\right)}{\left|S\right|}$. $G$ is $d-regular$ with eigenvalues $d=\lambda_{1}\ge\lambda_{2}\ge\dots\ge\lambda_{n}$.
\begin{insight}
The graph is connected iff $\lambda_{1}>\lambda_{2}$%
\end{insight}

\begin{thm}
\begin{insight}
page 42 in ``AllLectures''%
\end{insight}
Let $G$ be a $d-regular$ connected non-bipartite graph with eigenvalues
$d=\lambda_{1}\ge\lambda_{2}\ge\dots\ge\lambda_{n}$ then $\frac{d-\lambda_{2}}{2}\le h\left(G\right)\le\sqrt{d^{2}-\lambda_{2}}$.
\begin{proof}
we are only proving the bottom bound: $\lambda_{2}=\max\limits _{x\perp\one}\frac{xAx^{T}}{\norm x^{2}}$
\begin{insight}
$x\perp\one$ is beause the first eigenvector is $\one$ . and we
use the Rayleigh rits where $\lambda_{2}=\max\limits _{x\perp x_{1}}\frac{xAx^{T}}{\norm x^{2}}$%
\end{insight}
.We look at some partition of $V$: $S,\overline{S}$. What do we
know about the problem that we can use? We know that the graph is
$d-reg$ and we know the partition $S,\overline{S}$. So we will define
$x$ as follows:
\[
x_{i}=\begin{cases}
\left|\overline{S}\right| & i\in S\\
-\left|S\right| & i\in\overline{S}
\end{cases}
\]
this vector $x$ holds that $x\perp\one$ ($\left\langle x,\one\right\rangle =\sum x_{i}\cdot1=\left|\overline{S}\right|\cdot\left(-\left|S\right|\right)+\left|S\right|\cdot\left|\overline{S}\right|=0$)%
\begin{insight}
This is because we have $\left|\overline{S}\right|$ times $\left(-\left|S\right|\right)$
in entries in $x$ and $\left|S\right|$ times $\left|\overline{S}\right|$
entries in x.%
\end{insight}
. Now we need to calculate $\max\limits _{x\perp\one}\frac{xAx^{T}}{\norm x^{2}}$.
First: 
\begin{align*}
\norm x^{2} & =\sum_{i=1}^{n}x_{i}^{2}=\\
 & =\left|S\right|\left|\overline{S}\right|^{2}+\left|S\right|^{2}\left|\overline{S}\right|=\left|S\right|\left|\overline{S}\right|\left(\left|\overline{S}\right|+\left|S\right|\right)=n\left|S\right|\left|\overline{S}\right|
\end{align*}
 
\begin{align*}
xAx^{T} & =\sum a_{ij}x_{i}x_{j}=
\end{align*}
 this is like summing where $a_{ij}=1$ (where $i$ is a neighbour
of $j$) so 
\[
=\sum\limits _{i\sim j}x_{i}x_{j}=\underbrace{2\cdot e\left(S\right)\cdot\left|\overline{S}\right|^{2}}_{\left(1\right)}+\underbrace{2\cdot e\left(\overline{S}\right)\cdot\left|S\right|^{2}}_{\left(2\right)}-\underbrace{2e\left(S,\overline{S}\right)\left|S\right|\left|\overline{S}\right|}_{\left(3\right)}
\]
\end{proof}
\begin{enumerate}
\item $i,j$ both are vertices in $S$ hence $x_{i}\cdot x_{j}=\left|\overline{S}\right|\cdot\left|\overline{S}\right|$
and we count them $2\cdot e\left(S\right)$ times. 
\item same explanation as above but for $i,j$ vertices in $\overline{S}$.
\item $i\in S$, $j\in\overline{S}$, $x_{i}\cdot x_{j}=\left|\overline{S}\right|\cdot\left(-\left|\overline{S}\right|\right)$.
and we count them $2\cdot e\left(S,\overline{S}\right)$ times.
\end{enumerate}
Another important observation is $d\cdot\left|S\right|=2\cdot e\left(S\right)+e\left(S,\overline{S}\right)$
because this is a d-regular graph - we count for each vertex the number
of times it is touched by the the edges in $S$ plus the edges that
connect to$S$ from $\overline{S}$. (from the other side $d\left|\overline{S}\right|=2e\left(\overline{S}\right)+e\left(S,\overline{S}\right)$)
So %
\begin{insight}
\begin{align*}
xAx^{T} & =2\cdot e\left(S\right)\cdot\left|\overline{S}\right|^{2}+2\cdot e\left(\overline{S}\right)\cdot\left|S\right|^{2}-2e\left(S,\overline{S}\right)\left|S\right|\left|\overline{S}\right|=\\
* & =2\left(\frac{1}{2}\left(d\left|S\right|-e\left(S,\overline{S}\right)\right)\left|\overline{S}\right|^{2}+\frac{1}{2}\left(d\left|\overline{S}\right|-e\left(S,\overline{S}\right)\right)\left|S\right|-e\left(S,\overline{S}\right)\left|\overline{S}\right|\cdot\left|S\right|\right)=\\
* & =d\left|S\right|\left|\overline{S}\right|\left(\left|\overline{S}\right|+\left|S\right|\right)-e\left(S,\overline{S}\right)\left(\left|\overline{S}\right|^{2}+2\left|\overline{S}\right|\left|S\right|+\left|S\right|^{2}\right)=\\
 & =d\left|S\right|\left|\overline{S}\right|n-e\left(S,\overline{S}\right)n^{2}
\end{align*}
( {*} )were taken from allLectures.%
\end{insight}
{} 
\[
xAx^{T}=\dots=d\left|S\right|\left|\overline{S}\right|n-e\left(S,\overline{S}\right)n^{2}
\]

so $\lambda_{2}\ge\frac{xAx^{T}}{\norm x^{2}}$ lets assign what we
calculated: $\lambda_{2}\ge\frac{d\left|S\right|\left|\overline{S}\right|n-e\left(S,\overline{S}\right)n^{2}}{n\cdot\left|S\right|\left|\overline{S}\right|}$
since $\left|\overline{S}\right|\ge\frac{1}{2}n$ 
\[
\lambda_{2}\ge\frac{d\left|S\right|\left|\overline{S}\right|n-e\left(S,\overline{S}\right)n^{2}}{n\cdot\left|S\right|\left|\overline{S}\right|}\ge d-2\frac{e\left(S,\overline{S}\right)}{\left|S\right|}
\]
organize this a bit more to get: 
\[
\frac{e\left(S,\overline{S}\right)}{\left|S\right|}\ge\frac{d-\lambda_{2}}{2}
\]
We saw that $\forall S\,\;\left|S\right|\le\frac{n}{2}$ then $\frac{e\left(S,\overline{S}\right)}{\left|S\right|}\ge\frac{d-\lambda_{2}}{2}$.
This means that this is true for the minimum too $h\left(G\right)=\min\limits _{\begin{array}{c}
\emptyset\ne S\subseteq V\\
\left|S\right|\le\frac{\left|V\right|}{2}
\end{array}}\frac{e\left(S,\overline{S}\right)}{\left|S\right|}\ge\frac{d-\lambda_{2}}{2}$. 
\end{thm}

\begin{thm}
Fix $d$, let $G=G_{n}$ be an infinite family of $d-regular$ graphs.
then the following are equivalent:
\begin{enumerate}
\item $h\left(G_{n}\right)$ is bounded away from zero. %
\begin{insight}
    Bounded away from zero means that there exists a positve number $\epsilon>0$ that all elements are greater than.
\end{insight}
\item The spectral gap $\lambda_{1}-\lambda_{2}$ is bounded away from zero. 
\item The random walk on $G_{n}$ convergence in logarithmic time.
\end{enumerate}
\begin{insight}
In terms of apllication this is a huge thing. in all things ofenumaration
- checking different models this means that these things a re practical.
this means we don't have to wait.%
\end{insight}
\begin{proof}
\textcolor{red}{I Had a big problem taking notes in this part, I'll
try and fix this proof as soon as possible} ($2\then3$, What was
mostly intuition and not formal)$A$ is the matrix of the graph (a
stochastic matrix) so $A=\frac{1}{d}A_{G}$. We know that 
\[
\frac{A}{d}^{N}\to\text{\ensuremath{\begin{pmatrix}\\
 & \frac{1}{n}\\
\\
\end{pmatrix}}}
\]
Now assume we have some distribution, and we want to understand $\left(\frac{A_{G}}{d}\right)^{N}v\to\frac{1}{n}\one$.
We know nothing about $v$. We will express $v$ with the basis of
eigenvectors. $v=c_{1}\cdot\one+\sum c_{i}v_{i}$. where $\frac{A_{G}}{d}v_{i}=\lambda_{i}v_{i}$.
so what is $c_{1}$? $c_{1}=\left\langle v,\one\right\rangle $ and
because $v$ is a probability vector $c_{1}=\frac{1}{\sqrt{n}}$ .
thus $\frac{d}{\sqrt{n}}$ normalizes what happen when we apply $A^{N}$
to $v$? $A^{N}v=c_{1}\cdot A\one+\sum c_{i}Av_{i}$something that
goes exponentially fast to zero. And we stopped here.%
\begin{insight}
\begin{proof}
We are proving that $2\then3$.
\[
\left(\frac{A}{d}\right)^{n}v\to\frac{1}{n}\one
\]
we want to write $v$ as $v=c_{i}\frac{1}{\sqrt{n}}\one+\sum c_{i}v_{i}$.
where $c_{i}=\left\langle v,v_{i}\right\rangle $. Because $v$ is
a prob vector $c_{1}=\frac{1}{\sqrt{n}}..$ when $A^{n}v$ we get
that $\one$ stays the same, and all the other vectors other then
$\one$ tends to zero.
\end{proof}
Question given $d$ how large can the spectral be? $\iff$how small
can $\lambda_{2}$ be?%
\end{insight}
\end{proof}
\end{thm}

Question: given $d$ how large can the spectral gap be? $\iff$ how
small can $\lambda_{2}$ be?
\begin{thm}[Alon Boppana bound]
 %
\begin{insight}
Pg 109 in Ayelet's sum%
\end{insight}
$\lambda_{2}\ge2\sqrt{d-1}-o_{n}\left(1\right)$ .

a hint was given but not an actual proof.
\end{thm}

\begin{thm}
$\lambda_{2}\ge\sqrt{d}$
\begin{proof}
Let $A$ be the adjacency matrix of $G$%
\begin{insight}
G is d regular and have eigenvalues and stuff%
\end{insight}
. Then $tr\left(A^{2}\right)=\sum\lambda_{i}^{2}$, $\lambda_{1}=d$
so the first in the sum is $\lambda_{1}^{2}=d^{2}$ so $\sum_{i=1}\lambda_{i}=d^{2}+\sum_{i=2}\lambda_{i}^{2}$
. We know that $\lambda_{2}\ge\lambda_{3}\ge\lambda_{4}\dots\ge\lambda_{n}$
thus $\sum_{i=1}\lambda_{i}=d^{2}+\left(n-1\right)\cdot\lambda_{2}^{2}$
on the other hand 
\[
tr\left(A^{2}\right)=n\cdot d
\]
because $A^{2}$ in the $i,j$ entry counts how many paths of size
2 there are between $i$to $j$ ($d$ paths). so there are $n$ times
$d$ paths. so
\[
tr\left(A^{2}\right)=nd\le d^{2}+\left(n-1\right)\lambda_{2}^{2}
\]
 we conclude: $d-o\left(d\right)=\frac{nd-d^{2}}{n-1}\le\lambda_{2}^{2}$
\end{proof}
\end{thm}

